% PDF 165--173；原书第 136--144 页。
\section{指令调优}
\label{sec:ch04-instruction}
第~\ref{sec:ch03-prompt-finetuning}~节介绍了大语言模型基于提示的微调。基本流程是使用表示给定任务的有标签问答对数据集，设计提示模板，把问题转换成自然语言指令，再用这些指令微调模型，使其预测人工整理的答案。这个办法很适合完成指定任务。不过，ChatGPT 和 Claude 2 等当时表现先进的聊天机器人，除了擅长训练所针对的问题，还能将能力扩展到微调目标之外的多种主题。怎样调优模型，使它无需示例便能准确完成与训练数据不同的任务？一种方法是指令调优（instruction tuning，IT）。

\subsection{指令调优流程}
从最基本的意义上说，指令调优就是把提示写成模型可执行的自然语言指令，再用这些提示微调语言模型。提示通常包含三个部分，即描述待完成任务的指令、完成任务所需的上下文，以及引出答案的提示。仍以第~\ref{sec:ch03-prompt-finetuning}~节的例子为例。
\begin{quote}
\texttt{Cannot watch this movie. This is [z].}
\end{quote}

做基于提示的微调时，我们收集大量这样的电影评论句子，将“great”或“terrible”分别填入提示，表示正面或负面类别，再用填好的提示调优模型。也可以为分类任务设计另一种模板，给出明确的自然语言指令，使模型照做。这样便无需从完形填空式提示中推断隐含的指示。

\begin{lstlisting}[nolol=true]
## Instruction: Please read the following film review and
determine if it is a positive or negative review. Respond
'positive' or 'negative', according to your classification.
## Review: Cannot watch this movie.
## Response:
\end{lstlisting}

这种表达更明确，也更接近人们习惯的交流方式。在许多指令模板样例上微调后，大语言模型除了能完成电影评论填空，还能学会遵循自然语言指令，并按人们期望的方式生成回答。

这种办法的一个突出特点是，遵循指令的能力可以扩展到微调中没有明确训练过的新任务。指令调优使用以自然语言表述的问题和目标答案，微调大语言模型，覆盖阅读理解、情感分析、翻译、常识推理等多种自然语言任务。任务既可以是有唯一答案的离散任务，例如“\texttt{What is the capital of Canada?}”，也可以是没有唯一正确答案的开放任务，例如“\texttt{Write a poem about pirates:}”。对后者，同一提示可以提供多条示范。

利用指令模板，把这些训练集转换为自然语言指令，调优模型，使它预测与每个训练输入配对的答案。调优以词元为单位进行。先给出完整的提示，通过反向传播调优模型，使它预测回答的第一个词元。随后，给出提示和第一个词元，使模型预测第二个词元，如此继续下去。多项研究（\refcite{ch04-wei-2021}{Wei et al., 2021}；\refcite{ch04-ouyang}{Ouyang et al., 2022}；综述见 \refcite{ch04-zhang-2023c}{Zhang et al., 2023c}）显示，经过指令调优的大语言模型，在微调时未接触过的自然语言处理任务上表现有所改善，可以很好地泛化。

\begin{figure}[!htbp]
\centering
% 重绘自原书第 138 页的图 4.2。
\begin{tikzpicture}[
  font=\footnotesize\sffamily,
  stagebox/.style={draw=black!44,fill=cyan!20,align=center,inner sep=3pt,
    minimum width=31mm,minimum height=12mm,rounded corners=2.4pt},
  result/.style={stagebox,fill=green!23},
  arrow/.style={draw=black!80,-{Stealth[length=1.4mm]},line width=.65pt}]
\node[stagebox,fill=yellow!27,minimum width=24mm] (pre) at (0,0) {预训练\\语言模型};
\node[stagebox] (ptft) at (3.20,1.66) {针对任务 A\\微调};
\node[stagebox,fill=red!19] (prompt) at (3.20,0) {为任务 A 设计零样本\\或少样本提示};
\node[stagebox] (it) at (3.20,-1.66) {在任务 B、C、D、\dots\\上微调};
\foreach \item/\title in {ptft/预训练再微调,prompt/基于提示的学习,it/指令调优}{
  \node[above=1pt of \item] {\title};
  \node[result,right=5mm of \item] (\item-result) {在任务 A 上推理};
  \draw[arrow] (\item.east) -- (\item-result.west);
}
\draw[arrow] (pre.east) to[out=13,in=190] (ptft.west);
\draw[arrow] (pre.east) -- (prompt.west);
\draw[arrow] (pre.east) to[out=-13,in=170] (it.west);
\end{tikzpicture}

\caption{预训练再微调、基于提示的推理和指令调优的总体流程比较。预训练再微调针对单一任务微调，再在同一任务上推理。基于提示的学习设计一条提示，可以包含上下文示例，随后用该提示推理。指令调优在众多数据集和不同任务上微调，使模型的能力可以泛化到训练中未见过的新任务。}
\label{fig:ch04-workflows}
\end{figure}

指令调优与第~3.1.2~节介绍的预训练再微调，以及第~3.2.2~节详细讨论的基于提示的学习和推理密切相关。图~\ref{fig:ch04-workflows} 显示了这些方法之间的关系。

\begin{keypoint}
预训练再微调对大语言模型做监督微调，使其完成特定的机器学习任务，涵盖基于任务头的微调与基于提示的微调。指令调优也做监督微调，但用统一的模板组织来自许多不同自然语言推断任务的输入与输出对。将输入统一写成自然语言提示后，模型学会执行这些任务，也学会正确理解人所写指令的含义。自回归大语言模型原本依据训练数据预测字符串中最可能出现的下一个词元。要让它给出有用、准确的指令回答，这一过程十分重要。
\end{keypoint}

下面讨论研究者如何收集指令调优数据集，并列举文献中使用指令调优实现领域适配的例子。

\subsection{指令调优数据}
指令调优训练集由问答对组成，将不同自然语言处理任务表述成自然语言指令及所需回答。研究者采用过多种方法来构造数据集，主要有以下三类。

\begin{enumerate}
\item 用提示模板把已有的自然语言处理训练数据集转换为指令调优数据。
\item 收集人编写的提示，再由人编写答案。
\item 收集人编写的提示，用已有的大语言模型生成答案。
\end{enumerate}

为了增加训练集的多样性和规模，往往会结合使用这三种方法。下面简述每种方法及其操作方式。图~\ref{fig:ch04-instruction-data} 给出概览。

\subsubsection{用模板转换的数据集做指令调优}
Hugging Face Hub 和 TensorFlow Datasets 等来源提供了大量可用于指令调优的自然语言处理数据集。这些数据集的结构和用途多种多样，一般有一个或多个输入，例如问题、上下文和指令，以及一个或多个可能的目标输出，输出可能按偏好排序。某些任务有明确的是非答案，如情感的正负判断。另一些需要从提供上下文的段落中抽取信息，如开卷问答。还有些任务比较开放，没有唯一正确的答案，如摘要和翻译。任务和主题的多样性，使数据能够涵盖许多相关的自然语言处理任务及领域知识。

\sourcefigure{fig04_03.pdf}{构造指令调优数据集的三种办法。（1）收集不同推理任务的开源数据集，用统一模板将问答转换成指令与输出，再做监督微调。（2）收集大量指令，例如用户发送给 OpenAI API 的查询，请人工标注者编写一个或多个回答，再微调基础模型。（3）收集种子任务及问答对，用大语言模型生成新任务及问答对，筛去质量较低的数据，再用筛选后的问答对微调基础模型。}{fig:ch04-instruction-data}

要用这些数据进行指令调优，可为每个数据集设计模板，将数据从原有结构转成自然语言指令与目标答案。以清单~\ref{lst:ch04-boolq} 中原书取自 GLUE BoolQ 的上下文、问题、答案三元组为例。

\noindent\begin{minipage}{\linewidth}
\begin{lstlisting}[caption={原书的 GLUE BoolQ 样例},label={lst:ch04-boolq}]
passage: "Look What You Made Me Do" is a song recorded by
American singer -songwriter Taylor Swift, released on August
24, 2017 by Big Machine Records as the lead single from her
sixth studio album Reputation (2017). Swift wrote the
song with her producer Jack Antonoff....
question: "did taylor swift write look what you made me do"
label: 1 (yes)
\end{lstlisting}
\end{minipage}\par

数据集的每条记录包含提供上下文的段落、一个关于该段落的问题，以及一个是非答案。可以按图~\ref{fig:ch04-instruction-data} 最上方流程中的模板 1 将它整理成下面的形式。

\noindent\begin{minipage}{\linewidth}
\begin{lstlisting}[caption={问答样例},label={lst:ch04-qa}]
"Look What You Made Me Do" is a song recorded by American singer-
songwriter Taylor Swift, released on August 24, 2017 by Big
Machine Records as the lead single from her sixth studio
album Reputation (2017). Swift wrote the song with her
producer Jack Antonoff... .

Based on the paragraph above answer the following question: did
taylor swift write "look what you made me do ?"
yes
\end{lstlisting}
\end{minipage}\par

训练时，把问题作为初始条件交给模型，再微调模型，使其正确回答“yes”。同一个数据集可以用几个不同的提示模板。这有助于防止模型对单一模板的特定措辞过拟合，也能减轻单一选择可能存在的缺点。

主要通过整理已有数据进行指令调优的代表模型是 FLAN（\refcite{ch04-wei-2021}{Wei et al., 2021}），它由大语言模型 LaMDA（\refcite{ch04-thoppilan}{Thoppilan et al., 2022}）经指令调优得到。研究者从开源数据库收集 62 个有标签的自然语言处理训练数据集，按任务分成 12 组。其中大多数是阅读理解和情感分析等自然语言理解任务，少数是机器翻译和摘要等生成任务。他们为每个数据集设计多个模板。为了检验指令调优模型的泛化能力，每次留出一个任务组，在其他 11 组上调优，再按留出组上的表现评分。与未经微调的 LaMDA 相比，FLAN 的表现大幅提升，尤其是在翻译和闭卷问答任务上。它还优于 GPT-3 等规模大得多、未经指令调优的大语言模型。这表明，要充分发挥模型的能力，指令调优十分重要，无论模型的参数量如何。

\subsubsection{用人工生成的样例做指令调优}
另一种方法是编写或收集一组指令任务，再请人写出答案供模型学习。我们称它们为人工生成的样例。这种方法适合为写诗等自然语言生成任务收集数据。这些任务没有特定的正确答案，但某些输出比其他输出更符合人的偏好。用人工编写的文本训练，有助于模型在回答时采用更接近人的语言表达方式。人工样例的主要优点是，模型设计者可以更细致地选择问题主题和规定回答细节。人工标注耗时且昂贵，因此专门人工制作的数据集，往往比汇集已有数据集得到的数据集更小。

InstructGPT（\refcite{ch04-ouyang}{Ouyang et al., 2022}）是利用人工生成样例的重要模型。研究者收集发送给 OpenAI API 的查询，并补充一些人工编写的问题，以构成输入集。随后雇用一组标注者，为这些查询撰写答案，构成完整数据集。他们采用的训练方法与第~\ref{ch:alignment}~章讨论的基于人类反馈的强化学习密切相关。训练后，标注者往往更喜欢 InstructGPT 的输出。与作为参照的模型比较，InstructGPT 的胜率约为 73\%，FLAN 约为 30\%。这个结果一部分来自训练技术，也反映了微调用的输出由人编写。用人工编写的答案训练的模型，比用多个不同数据集模板化训练的模型更符合人类偏好。

\subsubsection{用大语言模型生成的样例做指令调优}
最后一种指令调优方法，是用大语言模型回答查询，生成数据集。可以用要训练的同一个模型，也可以用另一个具有希望通过微调获得的特性的模型。一般先让模型生成若干答案，再用程序或人工筛选出质量最好的回答。与人工作答相比，这种方法快得多，费用也低得多，但回答所能涵盖的范围受大语言模型能力限制。

如果把答案交回最初生成它们的模型训练，还会遇到明显的数据污染问题。这样能否带来新信息？我们有理由追问，模型用自己生成的回答微调，究竟能学到什么。一种解释是，先让模型为同一个问题生成许多答案，再选出最符合人类判断的一两个，用它们微调模型。这样，模型就更有可能输出人们认为好的答案。即使它不像使用人工生成数据微调那样接触到新信息，也能学到人们希望它表现出的倾向。

为用实验检验相近的想法，\refcite{ch04-wang-2023b}{Wang 等人（2023b）}开发了 Self-Instruct，一种迭代式自举框架。它接收一组种子任务及各任务的示例输入与输出，再借助 GPT-3 根据种子输入生成新的查询。随后把查询交给 GPT-3 作答，将查询与输出配对构造指令调优数据集。研究者使用多种质量筛选方法，确保回答的准确性和多样性。最终结果表明，在自然语言生成任务上，他们的模型明显优于未经调优的 GPT-3，也优于由更小的大语言模型 T5 经指令调优得到的模型。

\subsection{用指令调优做领域适配}
前面主要谈了怎样推理未见过的任务。指令调优也常用于让大语言模型适应特定领域。研究表明，当任务高度专门化，需要理解专业术语或掌握预训练数据以外的信息时，经过领域适配的指令调优模型可以优于通用聊天机器人。这里概述教育、医学和金融领域的例子，这些例子只是指令调优的一部分应用。

在教育领域，\refcite{ch04-zhang-2023d}{Zhang 等人（2023d）}发布了用于写作任务的 Writing-Alpaca。调优后，模型能够对写作提出改进建议，例如纠正语法、使表达更清楚、简化难懂的句子或改写文本。调优数据包含每项任务的一句简短指令、要改正的输入句子以及改正后的输出句子。例如：

\begin{lstlisting}[nolol=true]
###Instruction:
Fix grammatical errors in the text
###Input:
She went to the markt
###Response:
She went to the market
\end{lstlisting}

这样，模型学会发现并改正拼写错误以及其他语法错误。Zhang 等人用约 60\,000 个文本改进样例微调 Llama-7B，涵盖七种文稿编辑任务，并从 Stanford-Alpaca 项目（\refcite{ch04-taori}{Taori et al., 2023}）另外加入 52\,000 条通用指令提示。所得模型显著优于作为参照的基础模型以及用专门程度较低的数据集调优的模型，但仍未达到用数百万个文本编辑样例做预训练再微调的模型的表现。这表明，指令调优可以用规模相对有限的数据得到质量较好的模型，而达到最佳表现有时仍需要更大的数据集。

在医学领域，\refcite{ch04-li-2023b}{Li 等人（2023b）}提出了 ChatDoctor，也是对大语言模型 Llama-7B 进行指令调优的版本。研究者用真实的医患文字聊天记录训练模型，使它能针对患者提出的问题给出行动建议。这个项目用指令提示调优模型，依次完成三项任务。

\begin{enumerate}
\item 用指令提示教模型从患者的问题中提取关键词。
\item 用指令提示引导模型考虑从内部数据库与维基百科取得、与关键词相关的上下文，从中抽取与患者问题有关的信息。
\item 用指令提示让模型考虑抽取出的上下文，并向患者建议接下来可采取的措施。
\end{enumerate}

训练数据由程序从医患对话中抽取，用于构造约 100\,000 个按顺序组织的样例。经过指令调优的模型，在理解患者症状并提供与症状有关的建议方面有明显改善。这也说明，通过串联多条指令，可以完成单条提示难以处理的更复杂任务。医学领域还有经过指令调优的大语言模型，专门用于放射学和肿瘤学等更细分的学科。

金融领域也出现了为不同任务设计的指令调优模型。例如，\refcite{ch04-zhang-2023a}{Zhang 等人（2023a）}开发了 Instruct-FinGPT，自动判断财经标题表达的情感是正面、中性还是负面。

\sourcefigure{fig04_04.pdf}{用模板处理金融数据以调优 Instruct-FinGPT 的例子（\refcite{ch04-zhang-2023a}{Zhang et al., 2023a}）。左侧为句子及对应标签，中间为指令模板，右侧为指令提示及回答。}{fig:ch04-fingpt}

研究者选用文献中标有情感类别的财经标题，拟定十条不同的情感分类指令，再构造带有标注答案的提示。图~\ref{fig:ch04-fingpt} 给出了一些模板转换的例子。这些例子说明了领域适配的用处。许多财经标题使用专门的措辞和句法，需要模型更充分地学习这些特征，才能取得更好的结果。使用这些数据调优 Llama-7B 后，与作为参照的 Llama 模型和 ChatGPT 等当时表现先进的聊天机器人相比，$F_1$ 得分提高超过 20\%。这说明了以指令调优进行领域适配的价值。

以上几例说明了怎样组织数据并用它微调大语言模型，从而改善表现、拓展能力。后续各节将考察高效微调大语言模型的不同办法。
